\uuid{4lwd}
\titre{Majorant}
\auteur{bodin ; hanani ; mzari}
\organisation{exo7}
\licence{CC BY-NC-SA 4.0}
\datecreate{2020-01-20}
\niveau{L1}
\module{Analyse}
\chapitre{Propriétés de R}
\sousChapitre{Maximum, minimum, borne supérieure}
\difficulte{3}
\theme{Maximum, majorant, majorant}
\exo7qcmid{lille-1.reels.030}

\contenu{
\begin{question}
\qclassification{120.02}
\qtags{motcle=majorant}
On dit que $M \in \Rr$ est un majorant d'un ensemble $E \subset \Rr$ si pour tout $x\in E$, on a $x \le M$.
\begin{answers}
    \bad{Un intervalle non vide et différent de $\Rr$ admet toujours un majorant.}

    \good{Un intervalle non vide et borné admet au moins deux majorants.}

    \good{Un ensemble qui admet un majorant, en admet une infinité.}

    \bad{L'ensemble $\Nn$ admet une infinité de majorants.}
\end{answers}
\begin{explanations}
L'ensemble des majorants (s'il est non vide) est du type $[M,+\infty[$.
\end{explanations}
\end{question}
}
